%================================================================
% 作業ログ 2026-06-16（test_v3 品質向上・4台輪唱・机上検証・KiCad回路図）
%   test_v3 品質向上（PR #25 で状態遷移/同期/Processing UI の徹底デバッグ），
%   4台輪唱（node_05 新設＋サイクル窓終端処理），机上検証スクリプトで
%   20アサーション全PASS＝ファームにバグなし確認，KiCad で指揮者テスト基板の
%   回路図作成まで．
%
%   ビルド:
%     cd work/shiozawa/work-0616/作業ログ0616
%     docker run --rm -v "$(pwd):/workspace" -w /workspace \
%       ghcr.io/paperist/texlive-ja:debian latexmk 作業ログ0616.tex
%================================================================
\documentclass[uplatex,dvipdfmx,a4paper,11pt]{jsarticle}

% --- 余白 ---
\usepackage[top=25mm,bottom=25mm,left=25mm,right=25mm]{geometry}

% --- 表・箇条書き・図・色 ---
\usepackage{array}
\usepackage{tabularx}
\usepackage{booktabs}
\usepackage{enumitem}
\usepackage{graphicx}
\usepackage{url}
\usepackage[table]{xcolor}
\usepackage{amssymb}

% --- 段落間隔（Wordのような見た目に寄せる）---
\setlength{\parindent}{0pt}
\setlength{\parskip}{0.6em}

% --- セクション見出しを「番号. 太字」+ 章末に水平線を入れる ---
\usepackage{titlesec}
\titleformat{\section}
  {\normalfont\Large\bfseries}{\thesection.}{0.6em}{}
\titleformat{\subsection}
  {\normalfont\large\bfseries}{\thesubsection}{0.6em}{}
\titlespacing*{\section}{0pt}{1.2em}{0.6em}
\titlespacing*{\subsection}{0pt}{1.0em}{0.4em}

% --- 章末水平線（手で \sectionrule と書いて区切る）---
\newcommand{\sectionrule}{\par\medskip\noindent\rule{\linewidth}{0.6pt}\par\medskip}

% --- 箇条書きの記号を Word 風（・→○→数字）に揃える ---
\renewcommand{\labelitemi}{\textbullet}
\renewcommand{\labelitemii}{\textopenbullet}
\setlist[itemize]{leftmargin=1.6em,itemsep=0.1em,topsep=0.2em}
\setlist[enumerate]{leftmargin=2.0em,itemsep=0.1em,topsep=0.2em}

% \textopenbullet が無い環境向けのフォールバック
\providecommand{\textopenbullet}{\ensuremath{\circ}}

% --- 進捗ガントチャート用（グリッド表）---
\definecolor{barDone}{gray}{0.55} % 完了
\definecolor{barProg}{gray}{0.78} % 進行中
\definecolor{barPlan}{gray}{0.90} % 予定
\definecolor{nowCol}{gray}{0.30}  % 今週の見出し
\newcolumntype{G}{>{\centering\arraybackslash}p{2.0em}}
\newcommand{\bD}{\cellcolor{barDone}}
\newcommand{\bP}{\cellcolor{barProg}}
\newcommand{\bQ}{\cellcolor{barPlan}}

\begin{document}

%================================================================
% ヘッダ
%================================================================
{\LARGE\bfseries 作業ログ}\par\bigskip

報告者：塩澤 匠生

報告日：2026年6月16日

プロジェクト名：タクトーン~ IMUジェスチャーで奏でるArduinoオーケストラ~ \\

\sectionrule

%================================================================
% 1. 進捗サマリー
%================================================================
\section{進捗サマリー（要約）}

\begin{itemize}
  \item 全体進捗率：70\%（test\_v3 の品質向上と 4 台輪唱を実装完了，机上検証でファームにバグなし確認．
        テストフェーズに移行し，実機切り分けが次の一手）
  \item 進捗状況（今週の変化）：
    \begin{itemize}
      \item PR \#25 で test\_v3 の状態遷移・同期・Processing UI を徹底デバッグし品質向上
      \item 4 台輪唱を実装（node\_05 新設，CANON\_CYCLE\_BEATS=56，サイクル窓終端処理）
      \item 輪唱ロジックの机上検証スクリプト（Python）で拍 1〜120 を全数検証し 20 アサーション全 PASS
      \item メニューナビを重力推定 LPF＋窓積算の縦/横判定に改善，状態遷移デッドタイム 1000ms 追加
      \item KiCad で ESP32-S3-DevKitC-1 の指揮者テスト基板回路図を作成
    \end{itemize}
  \item 今週の主要成果：
    \begin{itemize}
      \item PR \#25（\texttt{shiozawa-test\_v3-polish}）：node\_01 の Menu/Result 状態 IMU 監視・
            Fallback 復帰先修正，楽器クロック同期スナップ復帰（指揮者リセット 1 パケット追従），
            Processing 全画面に接続ステータス/操作ガイド追加，指揮者 \texttt{-std=gnu++17} で警告 8 件解消
      \item 4 台輪唱：新音色 4 種（トランペット/ホルン/トロンボーン/チューバ），
            headRestBeats 0/8/16/24 で 8 拍ずつ入りをずらす設計
      \item 机上検証で「コードにバグなし」を確信．実機で輪唱が聞こえなかった原因は
            楽器ファームの書き込み取り違え（test\_v2 のまま）の可能性が最有力
    \end{itemize}
  \item 重大課題（影響度：中）：
    \begin{itemize}
      \item {}[実機で node\_03/04 の輪唱が聞こえない．机上検証でバグなしのため，
            楽器ファームの書き込み取り違え（3 台とも同じノード用で遅延ゼロ）が有力仮説．
            test\_v3 main を全楽器に書き込み直して切り分ける必要あり]
    \end{itemize}
\end{itemize}

\sectionrule

%================================================================
% 2. 進捗ガントチャート（計画と現在地）
%================================================================
\section{進捗ガントチャート（計画と現在地）}

チームのガントチャート（\texttt{meetings/ガントチャーver1ト.xlsx}）のフェーズ区分（110 基本設計／
120 詳細設計／200 実装／300 テスト／400 ストレッチ）に，現在の進み具合を重ねたもの．
「今」の縦位置は本報告時点（6 月第 3 週）．

\medskip
{\footnotesize
\setlength{\tabcolsep}{2pt}
\renewcommand{\arraystretch}{1.2}
\begin{tabular}{@{}p{8.5em}|*{12}{G}@{}}
 & & & & & & & & & & \multicolumn{1}{c}{$\downarrow$今} & & \\
作業項目（フェーズ）
 & \cellcolor{barPlan}\tiny 4/15 & \cellcolor{barPlan}\tiny 4/22 & \cellcolor{barPlan}\tiny 4/29
 & \cellcolor{barPlan}\tiny 5/6 & \cellcolor{barPlan}\tiny 5/13 & \cellcolor{barPlan}\tiny 5/20
 & \cellcolor{barPlan}\tiny 5/27 & \cellcolor{barPlan}\tiny 6/3 & \cellcolor{barPlan}\tiny 6/10
 & \cellcolor{nowCol}\tiny\textcolor{white}{6/17} & \cellcolor{barPlan}\tiny 6/24 & \cellcolor{barPlan}\tiny 7/8 \\
\midrule
企画・テーマ決定        & \bD & \bD &     &     &     &     &     &     &     &     &     &     \\
110 基本設計            &     &     & \bD & \bD &     &     &     &     &     &     &     &     \\
120 詳細設計            &     &     &     &     & \bD &     &     &     &     &     &     &     \\
\textbf{200 実装}       &     &     &     &     &     &     &     &     &     &     &     &     \\
\quad test\_v2 輪唱・同期 &     &     &     & \bD & \bD & \bD & \bD &     &     &     &     &     \\
\quad test\_v3 ゲーム(firmware) &  &     &     &     &     & \bD & \bD & \bD & \bD &     &     &     \\
\quad test\_v3 Processing &   &     &     &     &     &     & \bD & \bD & \bD &     &     &     \\
\textbf{300 テスト}     &     &     &     &     &     &     &     &     &     &     &     &     \\
\quad 単体・結合・システム &   &     &     &     &     &     &     &     & \bP & \bP & \bQ &     \\
\quad 評価（精度・遅延） &     &     &     &     &     &     &     &     &     & \bQ & \bQ &     \\
400 ストレッチ（強弱ほか） &  &     &     &     &     &     &     & \bQ & \bQ &     &     &     \\
発表会                  &     &     &     &     &     &     &     &     &     &     &     & \multicolumn{1}{c}{$\bigstar$} \\
\bottomrule
\end{tabular}
}

\medskip
凡例：\colorbox{barDone}{~~}~完了 \quad \colorbox{barProg}{~~}~進行中 \quad
\colorbox{barPlan}{~~}~予定 \quad $\bigstar$~発表会（7/8）

\medskip
※ test\_v3 firmware は 6/10 週の品質向上＋4 台輪唱で完了．Processing も PR \#25 で
UI 改善が完了し，残りは実機テストでの微調整のみ．
※ 300 テストフェーズは机上検証スクリプト（20 アサーション全 PASS）で開始．
実機切り分け（書き込み取り違え確認）が次のマイルストーン．

\sectionrule

%================================================================
% 3. 作業ログ（詳細トラッキング）
%================================================================
\section{作業ログ（詳細トラッキング）}

\begin{tabularx}{\linewidth}{@{}p{10em}p{6em}p{4em}p{3em}X@{}}
\toprule
作業項目 & 実施日 & 作業時間 & 進捗 & コメント \\
\midrule
品質向上 PR \#25 & 2026-06-10 & 25 分 & 90\% & 状態遷移修正（Menu/Result IMU 監視・Fallback 復帰先），楽器クロック同期スナップ復帰，Processing 接続ステータス/操作ガイド追加，\texttt{-std=gnu++17} 警告解消 \\
ジャイロ 2D 表示削除 & 2026-06-10 & 5 分 & 100\% & ジャイロ値が不正のため 2D 表示機能ごと除去 \\
メニューナビ改善 & 2026-06-10 & 10 分 & 90\% & 重力推定 LPF＋窓積算の縦/横判定に変更（dynAcc 直読みより安定） \\
デッドタイム追加 & 2026-06-10 & 10 分 & 90\% & 状態遷移直後のジェスチャ/拍検出を 1000ms 抑制（誤動作防止） \\
4 台輪唱実装 & 2026-06-10 & 20 分 & 90\% & node\_05 新設，新音色 4 種，CANON\_CYCLE\_BEATS=56，サイクル窓終端処理 \\
机上検証スクリプト & 2026-06-10 & 15 分 & 100\% & Python で applyPattern のサイクル窓ロジックを再現，拍 1〜120 で 20 アサーション全 PASS \\
作業メモ更新 & 2026-06-10 & 5 分 & 100\% & activeContext/progress を修正 5 タスク完了＋机上検証結果で更新 \\
KiCad 回路図作成 & 2026-06-16 & 30 分 & 80\% & ESP32-S3-DevKitC-1-N32R16V，10 ピンコネクタ（I2C），スイッチ 3 個，LED 回路．指揮者ノードのテスト基板 \\
\midrule
\multicolumn{2}{@{}l}{\textbf{合計}} & \textbf{120 分} & & \\
\bottomrule
\end{tabularx}

\sectionrule

%================================================================
% 4. 課題管理
%================================================================
\section{課題管理}

\begin{tabularx}{\linewidth}{@{}p{8em}Xp{2.5em}p{2.5em}Xp{3em}@{}}
\toprule
課題 & 発生原因 & 影響度 & 優先度 & 対策 & 状態 \\
\midrule
実機で輪唱が聞こえない & node\_03/04 が入りの遅延なしで発音．机上検証でファームにバグなし→書き込み取り違えが有力 & 高 & 高 & test\_v3 main を全楽器に書き込み直し，Processing 画面で partId/音色が 3 台で異なるか確認する & 対応中 \\
ゲームモードの打ち切り & GAME\_LENGTH\_BEATS\allowbreak{}=32 のため node\_04 は 16 拍しか弾けない（仕様上の制約） & 低 & 低 & ゲーム長の延長 or 採点区間の調整＝設計判断．自由演奏には影響なし & 保留 \\
UDP マルチキャストのパケロス & ACK・再送なし & 中 & 中 & 連送 4 発＋WiFi 復帰時再 join で対策済み．実機テストで効果確認する & 継続 \\
\bottomrule
\end{tabularx}

\medskip
※影響度＝プロジェクトへのインパクト，優先度＝対応の緊急性

\sectionrule

%================================================================
% 5. AI利用状況と検証
%================================================================
\section{AI利用状況と検証(再現性重視)}

\subsection{利用内容}

\begin{itemize}
  \item 利用目的：test\_v3 品質向上のレビュー，4 台輪唱のサイクル窓設計，机上検証スクリプトの作成
  \item 入力（プロンプト要旨）：
    \begin{itemize}
      \item {}[PR \#25 の状態遷移・同期ロジックの不具合箇所のレビューと改善案]
      \item {}[4 台で 8 拍ずつ入りをずらす輪唱のサイクル窓終端処理の設計]
      \item {}[applyPattern のサイクル窓ロジックを Python で再現し，全拍を検証するスクリプトの生成]
    \end{itemize}
  \item 出力概要：
    \begin{itemize}
      \item {}[Menu/Result 状態での IMU 監視不足，Fallback 復帰先の誤り，楽器側のスナップ復帰欠如を特定し修正案を提示]
      \item {}[CANON\_CYCLE\_BEATS=56（8 拍$\times$7 区間）で 4 声部を収容するサイクル窓設計]
      \item {}[Python スクリプト（\texttt{tools/canon\_sim/}）で拍 1〜120 の発音パターンを全数検証，20 アサーション全 PASS]
    \end{itemize}
\end{itemize}

\sectionrule

\subsection{検証方法・ファクトチェック}

\begin{itemize}
  \item 検証手法：
    \begin{itemize}
      \item {}[全 6 ノード（node\_01〜05 + devkitc）を \texttt{pio run} でビルドし SUCCESS 確認（警告 0）]
      \item {}[机上検証スクリプトで applyPattern の C++ ロジックを Python に移植し，拍ごとの発音ノード・休符を全数照合]
      \item {}[processing-java build 通過を確認]
    \end{itemize}
  \item 最終判断：
    \begin{itemize}
      \item 採用．机上検証で「ファームにバグなし」を確信．実機で輪唱が聞こえない原因は
            楽器ファームの書き込み取り違え（test\_v2 のまま or 3 台同一ノード用）が最有力．
            実機書き込み直しで切り分ける
    \end{itemize}
\end{itemize}

\sectionrule

%================================================================
% 6. 次回計画
%================================================================
\section{次回計画}

\begin{itemize}
  \item 目標：
    \begin{itemize}
      \item {}[test\_v3 main を全楽器（node\_02/03/04，可能なら node\_05）に書き込み直し，
            書き込み取り違えの切り分けを完了する]
      \item {}[Processing 画面で partId（0x02/0x03/0x04）と音色（トランペット/ホルン/トロンボーン）が
            3 台で異なることを確認する]
      \item {}[自由演奏 60 拍以上で入りが 8 拍ずつズレること，拍 33 以降の尻すぼみ→無音 8 拍→
            拍 57 再入を聴感確認する]
      \item {}[300 テストフェーズを本格化し，精度・遅延の定量評価を開始する]
    \end{itemize}
  \item 達成基準：
    \begin{itemize}
      \item {}[4 台の輪唱が実機で聴感的に成立する（入りのズレが明確に聞こえる）]
      \item {}[書き込み取り違えが確定 or 棄却され，次の対策が明確になっている]
    \end{itemize}
  \item 期限：
    \begin{itemize}
      \item {}[6/24 までに実機切り分け完了＋テスト評価を開始し，発表準備に着手]
    \end{itemize}
\end{itemize}

\sectionrule

%================================================================
% 7. その他
%================================================================
\section{その他}

\begin{itemize}
  \item {}[机上検証（20 アサーション全 PASS）により「コードにバグなし」を確信できたことで，
        デバッグの焦点が実機側（書き込み取り違え・ハードウェア）に絞られた．
        これは今後の切り分けを大幅に効率化する]
  \item {}[KiCad 回路図は \texttt{hardware/} 配下に作成済み（未コミット）．
        ESP32-S3-DevKitC-1-N32R16V をメインに，I2C コネクタ・スイッチ 3 個・LED 回路を配置した
        指揮者ノードのテスト用基板設計]
  \item {}[ゲームモードは追加機能（ストレッチ）であり，遅れても自由演奏のみで発表は成立する位置づけ]
\end{itemize}

\sectionrule

%================================================================
% 8. 付録
%================================================================
\section{付録}

\begin{itemize}
  \item 添付資料：
    \begin{itemize}
      \item ソースコード：\texttt{firmware/test\_v3/}（共通ライブラリ＋node\_01〜node\_05），\texttt{pc\_app/test\_v3/}
      \item 机上検証スクリプト：\texttt{tools/canon\_sim/}
      \item 回路図（未コミット）：\texttt{hardware/}
      \item GitHub リポジトリ：\url{https://github.com/takushio2525/2026-hackathon-team23}
      \item 関連コミット：\texttt{53437e8}，\texttt{4b480a6}，\texttt{9f9490c}，\texttt{b656884}，\texttt{80ec121}，\texttt{390e4ac}
    \end{itemize}
  \item 添付範囲：
    \begin{itemize}
      \item 差分（test\_v3 品質向上・4 台輪唱・机上検証の 1 週間分のスナップショット）
    \end{itemize}
\end{itemize}

\sectionrule

\end{document}
